%HOMEWORK template for M 325K
\documentclass{article}
\headheight=7pt         \topmargin=14pt
\textheight=574pt       \textwidth=445pt
\oddsidemargin=18pt     \evensidemargin=18pt

%Begin package import

\usepackage{amsmath, amssymb, amsthm}
\usepackage{mathtools}
\usepackage{pdfpages}
\usepackage{tikz-cd}

%End package import
%Begin custom commands

\newcommand{\C}{\mathbb{C}}
\newcommand{\R}{\mathbb{R}}
\newcommand{\Q}{\mathbb{Q}}
\newcommand{\Z}{\mathbb{Z}}
\newcommand{\N}{\mathbb{N}}
\newcommand{\F}{\mathbb{F}}
\newcommand{\abs}[1]{\left\lvert #1\right\rvert}
\newcommand{\RM}[1]{\mathrm{#1}}
\newcommand{\BF}[1]{\textbf{#1}}
\newcommand{\e}{\epsilon}
\newcommand{\ceil}[1]{\lceil{#1}\rceil}
\newcommand{\floor}[1]{\lfloor{#1}\rfloor}
\newcommand{\rarr}[0]{\rightarrow}
\newcommand{\IM}[1]{\text{Im}(#1)}
\newcommand{\Hom}[0]{\text{Hom}}
\newcommand{\Pow}[1]{\text{Pow}(#1)}
\newcommand{\angles}[1]{\langle #1 \rangle}

%End custom commands
%Begin custom theorems

\newtheorem{innercustomgeneric}{\customgenericname}
\providecommand{\customgenericname}{}
\newcommand{\newcustomtheorem}[2]{%
\newenvironment{#1}[1]
{%
\renewcommand\customgenericname{#2}%
\renewcommand\theinnercustomgeneric{##1}%
\innercustomgeneric%
}
{\endinnercustomgeneric}
}

\newcustomtheorem{THRM}{Theorem}
\newcustomtheorem{LEM}{Lemma}
\newcustomtheorem{EX}{Exercise}
\newcustomtheorem{COR}{Corollary}
\newcustomtheorem{PROB}{Problem}
\newcustomtheorem{claim}{Claim}

\newenvironment{solution}
{\renewcommand\qedsymbol{$\blacksquare$}\begin{proof}[Solution]}
{\end{proof}}

\newenvironment{definition}
{\renewcommand\qedsymbol{$\blacksquare$}\begin{proof}[Definition]}
{\end{proof}}

%End custom theorems
\begin{document}

\title{Fields 2}
\author{Arham Lodha}%put your name here
\date{March 26, 2024}%enter the due date here
\maketitle

\begin{PROB}{1}\label{Problem 1.}
\end{PROB}
\begin{proof}
    Consider rings $A \leq B \leq C$. Suppose $C$ is $f.g$ over $B$ as a $B-$module and $B$ $f.g$ over $A$ as a $A-$module. Thus $\exists b_1, \dots, b_n \in B \ni \forall b \in B$, $b$ is a A-linear combination of $\left\{ b_1, \dots, b_n \right\}$ or $b = \sum_{k = 1}^{n}a_k b_k$ where $a_k \in A$. Similarly, $\exists c_1, \dots, c_n \in C \ni \forall c \in C, c = \sum_{k = 1}^{n} g_k c_k$ where $g_k \in B$. Thus $c = \sum_{i = 1}^{n} c_i \sum_{j = 1} a_j b_j = \sum_{i = 1}^{n}  \sum_{j = 1} a_j (b_j c_i)$. Thus c is a A-linear combination of $\left\{ b_jc_i \right\}_{i,j}$ which is finite since $\left\{ b_j \right\}_j$ and $\left\{ c_i \right\}_i$ finite. Thus C is a f.g over A. 

    The generators of $C$ are the set $\left\{ b_j c_i \right\}_{i, j}$. 
    
    Suppose $A$ is Noetherian and $C$ is finitely generated $A$ module.
    
    \begin{enumerate}
        \item Case 1: Suppose $C = R^m$. Induct on $m$. 
        
        Base case: Suppose $m = 1$. Every submodule of $C$ is a ideal of $A$ and thus is finitely generated by the definition of Noetherian. 
        
        Inductive Hypothesis: Suppose for $m = k$, the theorem is true. 
        
        Inductive Step: Suppose $m = k + 1$. Consider the projection map $\varphi: C \to R^{k}$ given by dropping the last entry: $\varphi(\alpha_1, \dots, a_m) = (\alpha_1, \dots, \alpha_{m - 1})$. $\ker \varphi \cong  R$ which is finitely generated. Let $W$ be a submodule of $C$. $\varphi(W) \subseteq R^{m - 1}$ and thus $\varphi(W)$ is finitely generated by induction. $\ker \varphi \cap W$ is a submodule of $\ker \varphi$ and thus is a ideal and thus finitely generated. Thus $W$ is finitely generated.
        
        \item Case 2: The general case: Let $C$ be a general finitely generated $R$ module. Then there is a surjective map from $\varphi : F^m \to C$. Given a submodule $W$ of C, the correspondence theorem tells us $U = \varphi^{-1}(W)$ is a submodule of $R^m$ and thus finitely generated and thus $W$ is finitely generated.        
    \end{enumerate}

    Since $B$ is submodule of $C$ its finitely generated as A module. Thus $C$ must be finitely generated over $B$, because otherwise it would be infintily generated over $A$ which is a contradiction.     



    Consider rings $A \leq B \leq C$. Suppose $C$ is free with finite rank over $B$ and $B$ free with finite rank over $A$ as a $A-$module. Thus $\exists b_1, \dots, b_n \in B \ni \forall b \in B$, $b$ is a A-linear combination of $\left\{ b_1, \dots, b_n \right\}$ or $b = \sum_{k = 1}^{n}a_k b_k$ where $a_k \in A$. Similarly, $\exists c_1, \dots, c_n \in C \ni \forall c \in C, c = \sum_{k = 1}^{n} g_k c_k$ where $g_k \in B$. Thus $c = \sum_{i = 1}^{n} c_i \sum_{j = 1} a_j b_j = \sum_{i = 1}^{n}  \sum_{j = 1} a_j (b_j c_i)$. Thus c is a A-linear combination of $\left\{ b_jc_i \right\}_{i,j}$ which is finite since $\left\{ b_j \right\}_j$ and $\left\{ c_i \right\}_i$ finite. Thus $C$ is free over $A$ and $rank_A C = |\left\{ b_j c_i \right\}_{i, j}| = rank_a(B) rank_B(C)$    

    
\end{proof}

\bigskip

\begin{PROB}{2}\label{Problem 2.}
\end{PROB}
\begin{proof}
    $A \subset B \subset C$ with $C$ finite over $A$. Since $A$ is a feild its only ideal is $0$ and thus finitely generated and thus $A$ is Noetherian. Thus $B$ is finite dimensional over $A$. and $C$ is finite dimensional over B. Since $C \cong B^n$ and $B \cong A^m$ and thus are free modules, $[C : A] = [B : A][C : B]$              
\end{proof}

\section{Section I}
\begin{PROB}{1.1}\label{Problem 1.1.}
\end{PROB}
\begin{proof}
    
\end{proof}

\bigskip

\begin{PROB}{1.2}\label{Problem 1.2.}
\end{PROB}
\begin{proof}
    Let $F$ be a field, not of characteristic 2, let $x^2 + bx + c = 0$. Suppose $\exists \delta \in F \ni \delta^2 = b^2 - 4c$. 
    
    \begin{align*}
        x^2 + bx + c &= \frac{\delta^2 - 2b \delta + b^2}{4} + \frac{b \delta - b^2}{2} + c \\
        &= \frac{1}{4}(\delta^2 - 2b \delta + b^2 + 2bd - 2b^2 + 4c) \\
        &= \frac{1}{4}(b^2 - 4c + b^2 - 2b^2 - 4c) = 0
    \end{align*}
    
    Suppose, $\not \exists \delta \in F \ni \delta^2 =b^2 - 4c$, 
    
    Assume the contrary, $x^2 + bx + c = 0$ has a solution.  
    \begin{align*}
        x^2 + bx + c &= x^2 + bx + \frac{b^2}{4} + c - \frac{b^2}{4} = 0 \\
        (x - b)^2 &= \frac{b^2}{4} - c \\
        4(x - b)^2 &= b^2 - 4c = \delta \\
        (x - b)^2 &= \frac{\delta^2}{4}
    \end{align*}

    Since $\delta$ doesn't exists we have reached a contradiction. Thus $x^2 + bx + c$ has not solution.   
\end{proof}
\bigskip

\begin{PROB}{1.3}\label{Problem 1.3.}
\end{PROB}
\begin{proof}
    $\overline{\R} = \R$, thus subfields of $\C$ which are contained in $\R$ are not dense in $\C$.
    
    Suppose, $\mathbb{K} \subseteq \C$ is a proper subfield of $\C$ such that $\mathbb{K} \not \subset \R$. Thus $\exists z_0 \in \mathbb{K} \ni z_0 \not \in \R$ thus $z_0 = a + bi$. $\mathbb{K}$ cannot have the reals, otherwise all $\mathbb{K} = \mathbb{\C}$. But it must have the rationals and the rational multiples of $z_0$.      
    
    Let $z \in \C$, we can write $z = c + dz_0$ if $z = x + iy$, let $d = \frac{y}{b}$ and $c = x - \frac{y}{d}a$. If $r, s \in \Q$, then $r +sz_0 \in \mathbb{K}$. By choosing $r, s$ arbitrarily close to $a, b$ we can make $r + sz_0$ as close to $z$. Thus $\mathbb{K}$ dense in $\C$.    
\end{proof}

\section{Section II}

\begin{PROB}{2.1}\label{Problem 2.1.}
\end{PROB}
\begin{proof}
    $\alpha^3 - 3\alpha + 4 = 0 \implies \alpha^3 = 3\alpha - 4$. 
    
    Thus, 

    \begin{align*}
        1(\alpha^2 + \alpha + 1) &= \alpha^2 + \alpha + 1
    \end{align*}
    \begin{align*}
        \alpha(\alpha^2 + \alpha + 1) 
        &= \alpha^3 + \alpha^2 + \alpha \\
        &= \alpha^2 + 4\alpha - 4
    \end{align*}
    \begin{align*}
        \alpha^2(\alpha^2 + \alpha + a) &= \alpha(\alpha^2 + 4\alpha - 4) \\
        &= \alpha^3 + 4 \alpha^2 - 4\alpha \\
        &= 4\alpha^2 + 3\alpha - 4 - 4\alpha \\
        &= 4 \alpha^2 - \alpha - 4
    \end{align*}

    Thus,
    \begin{align*}
        1 = (c \alpha^2 + b \alpha + a)(\alpha^2 + \alpha + 1) &= 4c \alpha^2 - c\alpha - 4c + b \alpha^2 + 4b \alpha - 4b + a \alpha^2 + a \alpha + a \\
        &=  \alpha^2(a + b + 4c) + \alpha(a + 4b - c) + (a - 4b - 4c)
    \end{align*}

    Which makes a matrix equation like below

    \begin{equation*}
        \begin{bmatrix}
            1 & 1 & 4 \\
            1 & 4 & -1 \\
            1 & -4 & -4
        \end{bmatrix}\begin{bmatrix}
            a \\ b \\ c
        \end{bmatrix} = \begin{bmatrix}
            0 \\ 0 \\ 1
        \end{bmatrix}
    \end{equation*}

    $a = \frac{17}{49}$, $b = -\frac{5}{49}$, $c = -\frac{3}{49}$.  
\end{proof}

\bigskip

\begin{PROB}{2.2}\label{Problem 2.2.}
\end{PROB}
\begin{proof} 
    Let $f(x) = x^n + \sum_{m = 0}^{m}(-1)^{n - m} a_m x^m$. $\alpha$ is a root of $f$ in extension field $K$. 
    
    Thus, 

    \begin{align*}
        \alpha^n + \sum_{m = 1}^{m}(-1)^{n - m} a_m \alpha^m + a_0 &= 0 \\ 
        \alpha^n + \sum_{m = 1}^{m}(-1)^{n - m} a_m \alpha^m  &= -a_0 \\
        \alpha(\alpha^{n - 1} + \sum_{m = 1}^{m}(-1)^{n - m} a_m \alpha^{m - 1}) &= -a_0\\
        \frac{1}{\alpha} &= -\frac{1}{a_0}\left(\alpha^{n - 1} + \sum_{m = 1}^{m}(-1)^{n - m} a_m \alpha^{m - 1}\right)
    \end{align*}
\end{proof}


\bigskip

\begin{PROB}{2.3}\label{Problem 2.3.}
\end{PROB}
\begin{proof}
    Let $\alpha = \sqrt[3]{2}$  $\omega = e^{2 \pi i / 3}$ and $\beta := \omega \sqrt[3]{2}$. The minimal polynomial of $\alpha$ is $x^3 - 2$. This is also the minimal polynomial of $\beta$. $x^3 - 2$ is irreducible over $\Q$. Thus by Proposition 15.2.8, $\Q[\alpha] \cong \Q[\beta]$. 
    
    \[\begin{tikzcd}
        {\Q[x] \ (f)} &&& {\Q[\alpha]} \\
        \\
        &&& {\Q[\beta]}
        \arrow["\cong"{description}, hook, from=1-1, to=1-4]
        \arrow["\cong"{description}, hook, two heads, from=1-1, to=3-4]
        \arrow["{\cong }"{marking, allow upside down}, draw=none, from=1-4, to=3-4]
    \end{tikzcd}\]

    Thus since $x_1^2 + \dots + x_2^2 = -1$ has no solutions in $Q[\alpha]$ it has no solutions in $Q[\beta]$.   
\end{proof}

\section*{Section III}


\begin{PROB}{3.1}\label{Problem 3.1.}
\end{PROB}
\begin{proof}
    Let $F$ be a field and let $\deg F[\alpha] = \deg F[\alpha^2] = 5$. $F \subset F[\alpha^2] \subset F[\alpha]$. By the multiplication of degree, $5 = [F[\alpha] : F] = [F[\alpha^2] : F][F[\alpha^2] : F[\alpha]] = 5[F[\alpha^2] : F[\alpha]] \implies [F[\alpha^2] : F[\alpha]] = 1 \implies F[\alpha^2] = F[\alpha]$.     
\end{proof}

\bigskip

\begin{PROB}{3.2}\label{Problem 3.2.}
\end{PROB}
\begin{proof}
    Let $f(x) = x^4 + 3x + 3$. 
    By the eisenstein criterion, $x^4 + 3x + 3 \mod 3 = x^3$ and $9 \nmid 3$. Thus $x^4 + 3x + 3$ is irreducible in $\Q$. Let $r$ be a generic root of $f$. $[\Q[r] : \Q] = 4$. $[\Q[\sqrt[3]{2}] : \Q] = 3$. Assume the contrary, $r \in \Q[\sqrt[3]{2}]$. Thus $\Q[r] \subset \Q[\sqrt[3]{2}]$  By the multiplication of degree, $4 = [\Q[\sqrt[3]{2}] : \Q] = [\Q[r] : \Q][\Q[r] : \Q[\sqrt[3]{2}]] = 3[\Q[r] : \Q[\sqrt[3]{2}]]$, however $3 \nmid 4$, thus $[\Q[r] : \Q[\sqrt[3]{2}]]$ is not a positive integer. Thus $\Q[r] \not \subseteq \Q[\sqrt[3]{2}]$. Thus by contradiction, $r \not \in \Q[\sqrt[3]{2}]$ and since $r$ is a generic root of $f$ no root of $f$ is in $\Q[\sqrt[3]{2}]$ and thus $f$ is irreducible over $\Q[\sqrt[3]{2}]$.  
\end{proof}

\bigskip

Let $\zeta_n = e^{2 \pi i / n}$ for the upcoming problems. 

\begin{PROB}{3.3}\label{Problem 3.3.}
\end{PROB}
\begin{proof}
    Assume the contrary that $\zeta_5 \in \Q[\zeta_7]$. Thus $\Q[\zeta_5] \subseteq \Q[\zeta_7]$. $[\Q[\zeta_5] : \Q] = 4$ and $[\Q[\zeta_7] : \Q] = 6$ by the cyclotomic polynomial. By multiplication of degrees, $4$ must divide 6 which is a contradiction. Thus by contradiction, $\zeta_5 \not \in \Q[\zeta_7]$.      
\end{proof}

\bigskip

\begin{PROB}{3.4a}\label{Problem 3.4a.}

\end{PROB}
\begin{proof}
    $\zeta_4 = i$
    \begin{enumerate}
        \item The minimal polynomial is $x^2 + 1$. 
        \item Assume the contrary, $i \in Q[\zeta_3]$. Then, $Q \subseteq Q[i] \subseteq Q[\zeta_3]$. Since $\left[Q[\zeta_3] : Q\right] = 2$ and $\left[Q[i] : Q\right] = 2$. $Q[i] = Q[\zeta_3]$. Thus $\zeta_3 = \frac{1}{2} + i\sqrt{3}/2$ can be written as a $a + bi$ where $a, b \in \Q$. However $\sqrt{3} / 2$ is irrational thus no rational number $b = \sqrt{3} / 2$. Thus a contradiction occurs and $i \not \in Q[\zeta_3]$. Thus the $x^2 + 1$ is still minimal over $\Q[\zeta_3]$.            
    \end{enumerate}
\end{proof}

\begin{PROB}{3.4b}\label{Problem 3.4b.}

\end{PROB}
\begin{proof}
    $\zeta_6 = e^{2 \pi i / 6} = e^{i \pi/ 3}$
    \begin{enumerate}
        \item $\zeta_6$ is a root of $x^3 + 1 = (x^2 - x + 1)(x + 1)$. It is not a root of $x + 1$ but is a root of $x^2 - x + 1$. The only possible rational roots of $x^2 - x + 1$ are $\pm 1$ which are not roots of the poly and thus $x^2 - x + 1$ is irreducible and thus is the minimal polynomial of $\zeta_6$ in $\Q$.            
        \item $\frac{1}{2} + i\sqrt{3}/2 = \zeta_6 = - \zeta_3^2 =-(-\frac{1}{2} - i \sqrt{3}/2)$. Thus $\zeta_6$ is the root of the polynomial $x - \zeta_3^2$. Thus the minimal polynomial of $\zeta_6$ in $\Q[\zeta_3]$ is $x - \zeta_3^2$                    
    \end{enumerate}
\end{proof}

\begin{PROB}{3.4c}\label{Problem 3.4c.}
\end{PROB}
\begin{proof}
    $\zeta_8 = e^{\pi i / 4}$
    \begin{enumerate}
        \item $\zeta_8$ is a root to $x^4 + 1$ which has no real roots, as $\forall r \in \R, r^4 > 0$. Thus $x^4 + 1$ is irreducible in $\Q$ and thus the minimal polynomial. 
        \item The roots of $x^4 + 1$ are $\zeta_8, \zeta_8^3, \zeta_8^5, \zeta_8^7$ whose squares are $\pm i$. By 3.4a, $i \not \in \Q[\zeta_3] \implies \zeta_8, \zeta_8^3, \zeta_8^5, \zeta_8^7 \not \in \Q[\zeta_3]$, thus $x^4 + 1$ cannot be factored into a linear term. If $x^4 + 1$ is a product of two irreducible quadratic in $\Q[\zeta_3][X]$, it must break into two quadratic terms one of which must have $\zeta_8$ as a root.
        
        Thus $(x - \zeta_8)(x - \zeta_8^3) = x^2 - x(\zeta_8 - \zeta_8^3) + \zeta_8^4 =  x^2 - xi\sqrt{2} - 1$

        $(x - \zeta_8)(x - \zeta_8^5) = x^2 - x(\zeta_8 - \zeta_5) + \zeta_8^6 = x^2 - i$. $-i \not \in \Q[\zeta_3]$. Thus not a valid quadratic. 
        
        
        $(x - \zeta_8)(x - \zeta_8^7) = x^2 - x\sqrt{2} + 1$
        
        Assume the contrary, $\sqrt{2} \in \Q(\sqrt{3}, i)$. Thus $\sqrt{2} = a + bi + c\sqrt{3} + di\sqrt{3} = (a + c\sqrt{3}) + i(b + d\sqrt{3}) \implies b, d =0$ and $c \neq 0$ because $\sqrt{2} \not \in \Q$. Thus $2 = (a + c\sqrt{3})^2 = a^2 + 2ac + 3c^2 \implies 3c^2 = 2 - a^2 - 2ac \implies \sqrt{3} = \frac{2 - a^2 - 2ac}{c} \implies \sqrt{3} \in \Q$. However $\sqrt{3} \not \in \Q$, thus by contradiction $\sqrt{2} \not \in \Q(\sqrt{3}, i)$.


        Assume the contrary, $i\sqrt{2} \in \Q(\sqrt{3}, i)$. Thus $i\sqrt{2} = a + bi + c\sqrt{3} + di\sqrt{3} = (a + c\sqrt{3}) + i(b + d\sqrt{3}) \implies a, c =0$ and $b \neq 0$ because $\sqrt{2} \not \in \Q \implies i\sqrt{2} \not \in \Q[i]$. Thus $2 = -(b + d\sqrt{3})^2 = -b^2 - 2bd - 3d^2 \implies -3d^2 = 2 + a^2 + 2ac \implies \sqrt{3} = -\frac{2 + a^2 + 2ac}{c} \implies \sqrt{3} \in \Q$. However $\sqrt{3} \not \in \Q$, thus by contradiction $i\sqrt{2} \not \in \Q(\sqrt{3}, i)$.      

        
        $\Q(\zeta_3) \subseteq \Q(\sqrt{3}, i)$ but $\sqrt{2}$ and $i\sqrt{2}$ not in $\Q(\sqrt{3}, i)$, thus $\sqrt{2}, i\sqrt{2} \not \in \Q(\zeta_3)$
        
        
        Thus $x^4 + 1$ cannot be reduced and thus is the minimal polynomial for $\zeta_8$ 
    \end{enumerate}
\end{proof}

\begin{PROB}{3.4d}\label{Problem 3.4d.}
\end{PROB}
\begin{proof}
    $\zeta_9 = e^{2\pi i / 9}$
    \begin{enumerate}
        \item $\zeta_9$ is a root to $x^9 - 1 = (x^3 - 1)(x^6 + x^3 + 1)$. $\zeta_9$ not a root to $x^3 - 1$. $\zeta_9$ is a root to $f(x) = x^6 + x^3 + 1$. $f(x)$ is irreducible iff $f(x + 1)$ since $x \to x + 1$ is a automorphism of $\Q[x]$. 
        
        Thus,

        \begin{align*}
            f(x + 1) &= (x + 1)^6 + (x + 1)^3 + 1 \\
            &= \sum_{j = 0}^{6} {6 \choose j} x^{j} + \sum_{j = 0}^{3} {3 \choose j} x^{j} + 1 \\
            &= x^6 + 6x^5 + 15x^4 + 20x^3 + x^3 + 15x^2 + 3x^2 + 6x + 3x + 1 + 1 + 1 \\
            &= x^6 + 6x^5 + 15x^4 + 21x^3 + 18x^2 + 9x+ 3
        \end{align*}

        By the eisenstein criterion, $f(x + 1)$ is irreducible thus $f(x)$ is irreducible. 
        
        \item $\Q \subseteq \Q[\zeta_3] \subseteq \Q[\zeta_9]$. 
        $\zeta_9^3 = \zeta_3$. 
        $6 = [\Q[\zeta_9] : \Q] = [\Q[\zeta_9] : \Q[\zeta_3]][\Q[\zeta_3] : \Q] = 2 [\Q[\zeta_9] : \Q[\zeta_3]] \implies [\Q[\zeta_9] : \Q[\zeta_3]] = 3 \implies$ the minimal polynomial is $x^3 - \zeta_3$.   
    \end{enumerate}
\end{proof}

\begin{PROB}{3.4e}\label{Problem 3.4e.}
\end{PROB}
\begin{proof}
    $\zeta_{10} = e^{\pi i / 5}$
    \begin{enumerate}
        \item $\zeta_{10}$ is a root of $x^5 + 1 = (x - 1)(x^4 - x^3 + x^2 - x + 1)$. $\zeta_{10}$ is not a root of of $x - 1$ but is a root of $f(x) = x^4 - x^3 + x^2 - x + 1$. Thus $[\Q[\zeta_{10}] : \Q] \leq 4$. $\zeta_{10}^2 = \zeta_5$. Thus $\Q[\zeta_{5}] \subseteq \Q[\zeta_{10}]$. 
        
        $\zeta_5$ is the solution of the polynomial $f(x) = x^4 + x^3 + x^2 + 1$. 
        
        \begin{align*}
            f(x + 1) &= x^4 + 4x^3 + x^3 + 6x^2 + 3x^2 + x^2 + x + 4x + 3x + 2x + 5 \\
            &= x^4 + 5x^3 + 10x^2 + 10x + 5
        \end{align*}

        $f(x + 1)$ is irreducible by eisenstein criterion. Thus $f(x)$ is irreducible.   
        
        But $[\Q[\zeta_5] : \Q] = 4$. Since $[\Q[\zeta_{10}] : \Q] = 4[\Q[\zeta_{10}] : \Q[\zeta_5]] \leq 4 \implies [\Q[\zeta_{10}] : \Q] = 4$ and $\Q[\zeta_5] = \Q[\zeta_{10}]$      
        \item Since $\Q[\zeta_5] = \Q[\zeta_{10}]$, it suffices to show that the minimal polynomial for $\zeta_{5}$ is not reducible over $\zeta_3$. Assume the contrary, $f(x)$ is reducible over $\Q[\zeta_3] \implies \Q \subseteq \Q[\zeta_3] \subseteq \Q[\zeta_5]$. Since all the intermediate fields of $\Q[\zeta_5]$ can be made by taking all the $g$ such that $g | f$ and adjoining all the coefficients of $g$ to $\Q$. The only roots of $f$ are $\zeta_5, \dots, \zeta_5^4$ and $[\Q[\zeta_5] : \Q[\zeta_3]] = \frac{4}{2} = 2$. So we have to take all possible pairs to see if all the coefficients are $\Q[\zeta_3]$. 
        
        Thus $(x - \zeta_5^i)(x - \zeta_5^j) = x^2 + (\zeta_5^i + \zeta_5^j)x + \zeta_5^{i + j}$. 
        
        If $j = 5 - i$, we get $x^2 + 2x\cos(\frac{2 \pi i}{5}) + 1$. $2x\cos(\frac{2 \pi i}{5})  \not \in \Q[\zeta_3]$. 
        
        If $j \neq 5 - i$. we get $x^2 + 2x(\zeta_5^i + \zeta_5^j) + \zeta_5^{i + j}$. But $\zeta_5^{i + j}$ generates all of $\mu_5$. Thus we get $\Q[\zeta_5]$ back which is not in $\Q[\zeta_3]$ . 
        
        
        Thus $x^4 + x^3 + x^2 + x + 1$ is irreducible so $x^4 - x^3 + x^2 - x + 1$ is irreducible over $\Q[\zeta_3]$   
    \end{enumerate}
\end{proof}

\begin{PROB}{3.4f}\label{Problem 3.4f.}
\end{PROB}
\begin{proof}
    $\zeta_{12} = e^{\pi i / 6}$
    \begin{enumerate}
        \item $\zeta_{12}$ is a root of $x^6 + 1 = (x + 1)^2(x^4 + x^2 + 1)$. $\zeta_{12}$ is not a root of of $x - 1$ but is a root of $f(x) = x^4 + x^2 + 1$. $f(x)$ is irreducible because it is the 12th cyclotomic polynomial. Thus f is the minimal polynomial. 
        \item $\Q \subseteq \Q[\zeta_3] \subseteq \Q[\zeta_{12}]$. $\zeta_{12}^2 = e^{\pi i / 3} = -\zeta_3^2$. $[\Q[\zeta_{12} : \Q[\zeta_3]]] = \frac{[\Q[\zeta_{12} : \Q]}{[\Q[\zeta_3] : \Q]} = \frac{4}{2} = 2$. Thus the minimal polynomial is $x^2 + \zeta_3^2$  
    \end{enumerate}
\end{proof}

\bigskip

\begin{PROB}{3.5}\label{Problem 3.5.}
\end{PROB}
\begin{proof}
    By the cyclotomic polynomial, we know $[\Q[\zeta_n] : \Q] = \varphi(n)$, or the number of relative numbers less than it. The only numbers with $\varphi(n) \leq 3$ are $2, 3, 4, 6$.   
\end{proof}

\bigskip

\begin{PROB}{3.6}\label{Problem 3.6.}
\end{PROB}
\begin{proof}
    Let $a \in \Q$ such that it is not a square in $\Q$. $\Q \subseteq \Q[\sqrt{a}] \subseteq \Q[\sqrt[4]{a}]$ and $[\Q[\sqrt{a}] : \Q] = 2$. Assume the contrary, $\sqrt[4]{a} \in \Q[\sqrt{a} \implies \exists x, y \in \Q \ni \sqrt[4]{a} = x + y\sqrt{a}$. $\sqrt{a} = (x + y\sqrt{a})^2 = x^2 + 2xy\sqrt{a} + y^2a = (x^2 + y^2a) + 2xy\sqrt{a} \implies x^2 + y^2a = 0$. But since $x^2, y^2, a > 0$, $x = 0$, $y = 0$. Which is a contradiction. Thus by contradiction, $\sqrt[4]{a} \not \in \Q[\sqrt{a}]$. But $\sqrt[4]{a}$ minimal polynomial in $\Q[\sqrt{a}]$ is $x^2  - \sqrt{a}$. Thus $[\Q[\sqrt[4]{a}] : \Q[\sqrt{a}]] = 2$. Thus $[\Q[\sqrt[4]{a}] : \Q] = 2 * 2 = 4$.                
\end{proof}

\bigskip

\begin{PROB}{3.7a}\label{Problem 3.7a.}
\end{PROB}
\begin{proof}
    Let $\alpha = \sqrt[4]{-2}$. 
    $\Q \subseteq \Q[\sqrt{-2}] \subseteq \Q[\alpha]$. 
    Assume the contrary, $i \in Q[\alpha]$. Thus, $\Q \subseteq \Q[i] \subseteq \Q[\alpha]$. $\sqrt{-2} = i\sqrt{2} \not \in Q[i]$. Thus $\Q[i] \subset \Q[i, \sqrt{-2}] \subseteq \Q[\alpha]$. $[\Q[\alpha] : \Q[i]] = [\Q[\alpha] : \Q]/[\Q[i] : \Q] = \frac{4}{2} = 2$. Since $2$ is prime, $\Q[i, \sqrt{-2}] = \Q[\alpha]$. Thus $\alpha = a + b\sqrt{-2}$ for $a, b \in Q[i]$. $\alpha^2 = \sqrt{-2} = a^2 + 2ab\sqrt{-2} - 2b^2 = (a^2 - 2b^2) + 2ab\sqrt{-2}$. Thus $0 = a^2 - 2b^2 =  (a - \sqrt{2}b)(a + \sqrt{2}b)$ this means, $a = \sqrt{2}b$ or $a = -\sqrt{2}b$ either of which is impossible to satify if $a, b \in \Q[i]$. $a = b =0$ is also a valid solution but this implies $\sqrt{-2} = 0$ which is a contradiction. Thus by contradiction, $i \not \in Q[\alpha]$.                  
\end{proof}

\begin{PROB}{3.7b}\label{Problem 3.7b.}
\end{PROB}
\begin{proof}
    Let $\omega = \zeta_3$, $\alpha = \sqrt[3]{2}$, and $\beta = \sqrt[3]{5}$. It suffices to prove that $\beta \not \in \Q[\alpha, \omega]$ because, $\Q[\beta] \subset \Q[\alpha, \omega]$. Since $3$ and $2$ are coprime, 
    
    \[\begin{tikzcd}
        & {\Q[\alpha, \omega]} \\
        {\Q[\omega]} && {\Q[\alpha]} \\
        & \Q
        \arrow["3"', from=3-2, to=2-3]
        \arrow["2", from=3-2, to=2-1]
        \arrow["3", from=2-1, to=1-2]
        \arrow["2"', from=2-3, to=1-2]
        \arrow["6"{description}, from=3-2, to=1-2]
    \end{tikzcd}\]

    and 
    \[\begin{tikzcd}
        & {\Q[\beta, \omega]} \\
        {\Q[\omega]} && {\Q[\beta]} \\
        & \Q
        \arrow["3"', from=3-2, to=2-3]
        \arrow["2", from=3-2, to=2-1]
        \arrow["3", from=2-1, to=1-2]
        \arrow["2"', from=2-3, to=1-2]
        \arrow["6"{description}, from=3-2, to=1-2]
    \end{tikzcd}\]

    Assume the contrary, $\Q[\beta, \omega] = \Q[\alpha, \omega]$. $\Q[\alpha, \omega]$ is generated by two automorphisms $\sigma$ which interchanges $\omega$ and $\omega^2$ and $\gamma$ which sends $\alpha \to \omega \alpha$ and sends $\omega \alpha \to \omega^2 \alpha$ but fixes $\omega$ and $\omega^2$. $\gamma$ can't fix $\beta$ or it would fix the whole real subfield of $\Q[\alpha, \beta]$ and it already fixes $\omega$ and $\omega^2$, so it would be trivial. But it must send $\beta$ to another cube root of $5$. 
    
    Suppose, $\gamma(\beta) = \omega \beta$. Then $\gamma$ and $\sigma$ fix $\frac{\beta}{\alpha}$ so $\frac{\beta}{\alpha} \in \Q \implies \sqrt[3]{\frac{5}{2}} \in \Q$. Thus $\exists a, b \in \Z \ni \gcd(a, b) = 1 \ni \frac{a}{b} = \frac{\beta}{\alpha} \implies \frac{a^3}{b^3} = \frac{5}{2} \implies 2a^3 = 5b^3$ thus $5 | a \implies 5 | b$ which is a contradiction.    
    
    Suppose then $\gamma(\beta) = \omega ^2 \beta$. Then $\gamma(\beta^2) = \omega \beta^2$. Thus $\gamma$ and $\sigma$ both fix $\frac{\beta^2}{\alpha} \implies \frac{\beta^2}{\alpha} \in \Q$. By the same argument as above this is a contradiction.
    
    Thus $\beta \in \Q[\alpha, \omega]$ produces a contradiction. Thus by contradiction, $\beta \not \in \Q[\alpha, \omega] \implies \beta \not \in \Q[\alpha]$.
\end{proof}


\bigskip

\begin{PROB}{3.8}\label{Problem 3.8.}
\end{PROB}
\begin{proof}
    Let $\alpha, \beta \in \C$. Suppose $\alpha + \beta$ and $\alpha \beta$ are algebraic numbers. Thus $\Q[\alpha + \beta, \alpha \beta]$ is finitely dimensional wrt to $\Q$ . Take $(x - \alpha)(x - \beta) = x^2 - (\alpha + \beta) + \alpha \beta \in Q[\alpha + \beta, \alpha \beta] [X]$. $\alpha$ and $\beta$ are roots of the polynomial. Thus $Q[\alpha , \beta]$ is a finite dimensional with respect to $\Q[\alpha + \beta, \alpha \beta]$. Thus $\Q[\alpha, \beta]$ is finite dimensional with respect to $\Q$. Since $\Q[\alpha] \subseteq \Q[\alpha, \beta]$ ad $\Q[\beta] \subseteq \Q[\alpha, \beta]$ they are also finite dimensional wrt to $\Q$ and thus $\alpha$ and $\beta$ are algebraic.       
\end{proof}

\bigskip

\begin{PROB}{3.9}\label{Problem 3.9.}
\end{PROB}
\begin{proof}
    Let $K = \Q[\alpha]$ and $L = \Q[\beta]$.  
    Let $m = \deg f$ and $n = \deg g$.      
    Suppose $f(x)$ is irreducible over $L[x]$. It is also irreducible over $\Q$. Thus $K \cong \Q[X]/f(x) \implies [K, \Q] = m$. Similarly, $\Q[\alpha, \beta] \cong L[X] / f(x) \implies [\Q[\alpha, \beta] : L] = m$. $L \cong \Q [x]/ g$ since $g$ is a irreducible polynomial where $\beta$. Thus $[L : Q] = n$. Thus since $[Q(\alpha, \beta) : \Q] = [L : \Q][\Q(\alpha, \beta) : L] = mn$. Thus $[\Q[\alpha, \beta] : K] = n \implies g$ is irreducible over $K[X]$.             
\end{proof}

\bigskip

\begin{PROB}{3.10}\label{Problem 3.10.}
    $K / F$ and $L / K$ algebraic $\implies$ $L / F$ algebraic    
\end{PROB}
\begin{proof}
    $\forall l \in L$, $\exists g \in K[X] \ni g(l) = 0$. Let $g(x) = \sum_{j = 0}^{n} k_j x^j$ where $k_j \in K$. 
    
    $[F[l, k_0, \dots, k_n] : F] = [F[l, k_0, \dots, k_n] : F[k_0, \dots, k_n]][F[k_0, \dots, k_n] : F]$
    
    $[F[l, k_0, \dots, k_n] : F[k_0, \dots, k_n]]$ is finite because $g \in F[k_0, \dots, k_n][X]$ and $g(l) = 0$. $[F[k_0, \dots, k_n] : F]$ is finite because $k_j$ is algebriac thus $\exists h_j \in F[X] \ni h_j(k_j) = 0$. Let $h(x_1, \dots, x_n) = \sum_{i = 0}^{n}h_i(x_i)$ and $(k_0, \dots, k_n)$ is a root of $h$. 
    
    Thus $[F[l, k_0, \dots, k_n] : F]$ is finite. Because $F \subseteq F[l] \subseteq [F[l, k_0, \dots, k_n]$,  $[F[l] : F] \leq [F[l, k_0, \dots, k_n] : F] \implies [F[l] : F] $ is finite and thus $l$ is algebraic over F. 
    
    Thus $L / F$ is algebraic.  
\end{proof}
\end{document}